% generated by roam.py handbook export -- не редактировать, править заметки
% 20260923143056-метопт_лекция_2_μ_сильно_выпуклые_и_l_гладкие_функции.org

\chapter{μ-сильно выпуклые и L-гладкие функции}\label{lecture:2}
\InputIfFileExists{text/02-intro}{}{}

\section{μ-сильно выпуклые функции}\label{lecture:2:1}  % конспект 2.1
\InputIfFileExists{text/02-1-intro}{}{}

\InputIfFileExists{text/53558BAA-700E-46BE-9EAB-6A2F65A5B4C0-before}{}{}
% generated by roam.py handbook export -- не редактировать, править заметки
% 20260922140145-неравенство_иенсена.org (body 8aec188530db)

\begin{theorem}[{Неравенство Йенсена}]\label[theorem]{note:53558BAA-700E-46BE-9EAB-6A2F65A5B4C0}
  Пусть $f$~--- \hyperref[note:1CA14F2C-637F-41DE-9E97-72075D19E97F]{выпуклая} на $E$, тогда
  \[
  \forall x_i \in E, 0 \le \lambda_i, \sum\limits_{i=1}^n \lambda_i = 1\quad f\left(\sum\limits_{i=1}^n \lambda_i x_i\right) \le
  \sum\limits_{i=1}^n \lambda_i f(x_i)
  \]
\end{theorem}

\begin{proof}
  Идея аналогична доказательству: \hyperref[note:42DFF6DC-CC9D-40AA-9D6F-218DCF613833]{Характеристика выпуклого множества через выпуклые комбинации}

  Доказывать будем по индукции.
  База: $n=1$,  тогда $\lambda_1 = 1 \Rightarrow f(\lambda_1 x_1) \le \lambda_1 f(x_1) \Leftrightarrow f(x_1) \le f(x_1)$. Очевидно;

  Давайте рассмотрим переход: $n \rightarrow n+1$ \\
  Имеем $\lambda_i \ge 0,\ \sum\limits_{i=1}^{n+1} \lambda_i = 1$. Если $\lambda_{n+1} = 1$, то мы опять вернемся к вырожденному случаю из базы,
  так как остальные $\lambda_i = 0, i < n+1$. Далее будем считать, что $\lambda_{n+1} \ne 1$.

  Пусть $s:= 1 - \lambda_{n+1}$, тогда
  \[
  \sum\limits_{i=1}^{n+1} \lambda_i x_i = \sum\limits_{i=1}^n \lambda_i x_i + \lambda_{n+1} x_{n+1} =
  s \cdot \frac{1}{s} \cdot \sum\limits_{i=1}^n \lambda_i x_i + \lambda_{n+1} x_{n+1}
  \]

  Рассмотрим $\frac{1}{s} \sum\limits_{i=1}^n \lambda_i x_i$. Это комбинация из $n$ точек, рассмотрим коэффициенты
  $\sum\limits_{i=1}^n \frac{\lambda_i}{s} = \frac{1-\lambda_{n+1}}{1-\lambda_{n+1}} = 1$, причем они все неотрицательные, так как $\lambda_i \ge 0$ и $s > 0$.

  Так как \hyperref[note:42DFF6DC-CC9D-40AA-9D6F-218DCF613833]{выпуклое множество содержит все свои выпуклые комбинации}, то $y := \sum\limits_{i=1}^n \frac{\lambda_i}{s} x_i \in E$.

  Итого, мы получили, что
  \[
  \sum\limits_{i=1}^{n+1} \lambda_i x_i = s y + \lambda_{n+1}x_{n+1} = (1-\lambda_{n+1}) y + \lambda_{n+1}x_{n+1}
  \]

  Воспользуемся определением \hyperref[note:1CA14F2C-637F-41DE-9E97-72075D19E97F]{выпуклости функции}.
  \[
  f( (1 - \lambda_{n+1})y + \lambda_{n+1} x_{n+1} ) \le (1-\lambda_{n+1}) f(y) + \lambda_{n+1} f(x_{n+1})
  \]
  Давайте теперь рассмотрим $f(y)$, это функция от \hyperref[note:8583F986-A6E8-4BBB-B2FB-BA0E9BBE58BF]{выпуклой комбинации} из $n$ точек.
  Значит мы можем применить индукционное предположение:
  \[
  f(y) = f \left(  \sum\limits_{i=1}^n \frac{\lambda_i}{s} x_i \right) \le \sum\limits_{i=1}^n \frac{\lambda_i}{s} f(x_i)
  \]

  Но тогда $(1-\lambda_{n+1}) f(y) = s f(y) \le \sum\limits_{i=1}^n \lambda_i f(x_i)$ \\
  Итого
  \[
  f \left( \sum\limits_{i=1}^{n+1} \lambda_i x_i  \right) = f(sy + \lambda_{n+1} x_{n+1}) \le sf(y) + \lambda_{n+1} f(x_{n+1}) \le \sum\limits_{i=1}^n \lambda_i f(x_i) + \lambda_{n+1} f(x_{n+1}) = \sum\limits_{i=1}^{n+1} \lambda_i f(x_i)
  \]
\end{proof}
  % Неравенство Йенсена -- нет в конспекте, открывало лекцию

\InputIfFileExists{text/22E546FC-186F-4A67-85AA-2C85DCEF9021-before}{}{}
% generated by roam.py handbook export -- не редактировать, править заметки
% 20260922143605-сильно_выпуклая_функция.org (body 7344ef77c2a1)

\begin{definition}[{Сильно выпуклая функция}]\label[definition]{note:22E546FC-186F-4A67-85AA-2C85DCEF9021}
  Пусть $E \subset \R^n$~--- \hyperref[note:316B5BFA-8B9F-4267-9086-D6593AAB5F9B]{выпуклое множество}, $f: E \rightarrow \R$.
  $f$~--- сильно выпуклая $\Leftrightarrow \exists \mu > 0: f(x) - \frac{\mu}{2}\|x\|^2$~--- \hyperref[note:1CA14F2C-637F-41DE-9E97-72075D19E97F]{выпуклая} \\
  Еще $f$ можно назвать $\mu$-сильно выпуклой.
\end{definition}

\begin{remark}
  Равносильное определение:
  \[
  \forall x, y \in E\ \forall \lambda \in [0;1]\ f(\lambda x + (1-\lambda)y) \le \lambda f(x) + (1-\lambda) f(y) - \frac{\mu}{2}\lambda(1-\lambda)\|x-y\|^2
  \]
\end{remark}
  % Сильно выпуклая функция -- опр. 2.1

\InputIfFileExists{text/1B5A3E25-56E8-4D5E-9CF6-E2731DA77604-before}{}{}
% generated by roam.py handbook export -- не редактировать, править заметки
% 20260922143839-дифференциальныи_критерии_сильнои_выпуклости.org (body 8cb481c501c7)

\begin{theorem}[{Дифференциальный критерий сильной выпуклости}]\label[theorem]{note:1B5A3E25-56E8-4D5E-9CF6-E2731DA77604}
  Пусть $E \subset \R^n$~--- \hyperref[note:316B5BFA-8B9F-4267-9086-D6593AAB5F9B]{выпуклое множество}. Если $f: E \rightarrow \R$~--- \hyperref[note:957FA09D-8883-40A5-B4C7-237D9F815A42]{дифференцируемая}, то
  $f$~--- $\mu$-\hyperref[note:22E546FC-186F-4A67-85AA-2C85DCEF9021]{сильно выпуклая} $\Leftrightarrow \forall x, y \in E\ f(y) \ge f(x) + \left\langle \nabla f(x); y-x \right\rangle + \frac{\mu}{2} \|y-x\|^2$
\end{theorem}

\begin{remark}
  Если для \hyperref[note:1CA14F2C-637F-41DE-9E97-72075D19E97F]{обычной выпуклости} мы подпирали функцию прямой (касательной), то в этом случае мы подпираем уже параболой
\end{remark}

\begin{remark}
  Если $\|y\| \rightarrow +\infty$, то правая часть тоже $\rightarrow +\infty$
\end{remark}

\begin{proof}
  Возьмем функцию $\phi(x) = f(x) - \frac{\mu}{2} \|x\|^2$. Она  дифференцируема.
  Докажем, что она выпуклая, используем \hyperref[note:18BD8497-CEBD-4962-8144-474B489606C3]{дифференциальный критерий выпуклости}.

  $\phi$~--- выпуклая $\Leftrightarrow \forall x,y \in E\ \phi(y) \ge \phi(x) + \left\langle \nabla \phi(x); y-x \right\rangle$

  Заметим $\nabla \phi(x) = \nabla f(x) - \mu x$, так как \hyperref[note:1756907B-1FF2-4F62-8553-759B8509281C]{градиент}~--- линеен, а градиент $\|x\|^2$~--- это $2x$.
  Давайте теперь раскроем $\phi$ в неравенстве:
  \[
  f(y) - \frac{\mu}{2} \| y \|^2 \ge f(x) - \frac{\mu}{2} \| x \|^2 + \left\langle \nabla f(x) - \mu x; y-x \right\rangle
  \]

  Давайте раскроем \hyperref[note:F2942E98-6115-4600-AA3F-E0F687D9C732]{скалярное произведение} в правой части:
  \[
  \left\langle \nabla f(x) - \mu x; y-x \right\rangle =  \left\langle \nabla f(x); y-x \right\rangle - \mu \left\langle x; y - x \right\rangle =  \left\langle \nabla f(x); y-x \right\rangle - \mu \left\langle x; y \right\rangle  +  \mu \left\langle x; x \right\rangle
  \]

  Тогда
  \[
  f(y) - \frac{\mu}{2} \| y \|^2  \ge f(x) - \frac{\mu}{2} \|x\|^2 + \left\langle \nabla f(x); y-x \right\rangle - \mu \left\langle x; y \right\rangle  +  \mu \left\langle x; x \right\rangle
  \]
  Перенесем норму $y$ на право и заменим $\mu \left\langle x;x \right\rangle = \mu \|x \|^{2}$,
  \[
  f(y) \ge f(x) - \frac{\mu}{2} \| x \|^2 + \mu \| x \|^2 + \frac{\mu}{2} \| y \|^2 + \left\langle \nabla f(x); y-x \right\rangle - \mu \left\langle x;y \right\rangle
  \]
  Итого получим
  \[
  f (y) \ge f(x) + \frac{\mu}{2} \| x \|^2 + \frac{\mu}{2} \| y \|^2 - \mu \left\langle x;y \right\rangle + \left\langle \nabla f(x); y-x \right\rangle
  \]

  Еще заметим, что  $\|y-x\|^2 = \|x\|^2 +\|y\|^2 - 2 \left\langle x;y \right\rangle$, давайте используем эту формулу \\
  Итого
  \[
  f(y) \ge f(x) + \frac{\mu}{2} \|y-x\|^2 + \left\langle \nabla f(x); y-x \right\rangle
  \]
  Все переходы~--- равносильности, поэтому получили в точности неравенство из условия: оба направления доказаны
\end{proof}
  % Дифференциальный критерий сильной выпуклости -- теорема 2.1 (в конспекте «критерий в терминах касательных»), замечание 2.1 — в /Rem/

\InputIfFileExists{text/51CC880D-BF19-43E2-9599-FB31A9A7367F-before}{}{}
% generated by roam.py handbook export -- не редактировать, править заметки
% 20260922144651-единственность_минимума_сильно_выпуклои_функции.org (body d3f1c9a34490)

\begin{theorem}[{Единственность минимума сильно выпуклой функции}]\label[theorem]{note:51CC880D-BF19-43E2-9599-FB31A9A7367F}
  Если $f: \R^n \rightarrow \R$~--- \hyperref[note:957FA09D-8883-40A5-B4C7-237D9F815A42]{дифференцируемая} и \hyperref[note:22E546FC-186F-4A67-85AA-2C85DCEF9021]{сильно выпуклая} функция, то у нее существует единственная точка минимума
\end{theorem}

\begin{remark}
  Дифференцируемость нужна только для существования: по \hyperref[note:1B5A3E25-56E8-4D5E-9CF6-E2731DA77604]{дифференциальному критерию} $f$ растет не медленнее параболы, поэтому минимум достигается на компактном шаре (Вейерштрасс). Единственность верна и без нее. Существование без дифференцируемости тоже верно, но доказывается сложнее~--- через непрерывность выпуклой функции на $\R^n$.
  Важно и то, что $f$ задана на всем $\R^n$: например, $f(x) = x^2$ на $(1;2)$ сильно выпукла, но минимум не достигается.
\end{remark}

\begin{proof}
  \proofpart{Существование}
  Давайте возьмем \hyperref[note:1B5A3E25-56E8-4D5E-9CF6-E2731DA77604]{дифференциальный критерий сильно выпуклой функции} для $x=0$
  \[
  f(y) \ge f(0) + \left\langle \nabla f(0); y \right\rangle + \frac{\mu}{2} \| y \|^2
  \]
  Применим \hyperref[note:D0601ACD-8C91-4A93-B0E4-AC0D529C6A5C]{КБШ} $| \left\langle x;y \right\rangle | \le \|x\| \cdot \|y\| \Leftrightarrow - \|x\| \cdot \|y\| \le   \left\langle x;y \right\rangle  \le \|x\| \cdot \|y\|$ (левую часть):
  \[
  f(y) \ge f(0) + \left\langle \nabla f(0); y\right\rangle + \frac{\mu}{2} \| y \|^2 \ge f(0) - \| \nabla f(0) \| \| y\| + \frac{\mu}{2} \| y \|^2
  \]

  Возьмем радиус $R := \frac{2 \| \nabla f(0) \|}{\mu}$, тогда, при $\|y\| > R$
  \begin{multline*}
  f(y) \ge f (0) -  \| \nabla f(0) \| \| y \| + \frac{\mu}{2} \| y \|^2 > f(0) - \| \nabla f(0) \| R + \frac{\mu}{2} R^2 = \\
   f(0) - \| \nabla f(0)\| \cdot  \frac{ 2 \| \nabla f(0)\|}{\mu} + \frac{\mu}{2} \cdot  \left( \frac{2 \| \nabla f(0) \|}{\mu} \right)^2 = \\
   f(0) - \frac{2}{\mu} \|\nabla f(0) \|^2 + \frac{2}{\mu} \| \nabla f(0) \|^2 = f(0)
  \end{multline*}
  Второе неравенство верно, так как парабола $h(r) = - \| \nabla f(0) \| r + \frac{\mu}{2} r^2$ возрастает при $r \ge \frac{\| \nabla f(0) \|}{\mu}$, а $\| y \| > R \ge \frac{\| \nabla f(0) \|}{\mu}$.
  Получили, что $f(y) > f(0)$, при $\| y \| > R$.

  Теперь рассмотрим $\| y \| \le R$~--- это замкнутый шар.
  Мы знаем, что $f$~--- дифференцируема, следовательно она непрерывна.
  Тогда по \hyperref[note:82C70263-336A-49F5-82DB-9D25C84D6A6D]{теореме Вейерштрасса} в $\overline{B}_R(0)$ у $f$ есть точка минимума на шаре $x^{*}$.
  Тогда  $\| y \| \le R \Rightarrow f(y) \ge f(x^{*}) \Rightarrow f(0) \ge f(x^{*})$.

  Тогда получается, что внутри шара это есть минимум, а вне шара $f(y) > f(0) \ge f(x^{*})$.
  Тогда $\forall y \in \R^n\ f(y) \ge f(x^{*})$.
  Тогда $x^{*}$~--- есть глобальный минимум

  \proofpart{Единственность}
  Пусть есть две точки минимума $x_1 \ne x_2$, $f(x_1) = f(x_2) =: m$.
  По определению \hyperref[note:22E546FC-186F-4A67-85AA-2C85DCEF9021]{сильной выпуклости}, $\phi(x) = f(x) - \frac{\mu}{2} \| x \|^2$  -- \hyperref[note:1CA14F2C-637F-41DE-9E97-72075D19E97F]{выпуклая}.
  По выпуклости $\phi$, $\phi(\lambda x_1 + (1-\lambda)x_2) \le \lambda \phi(x_1) + (1-\lambda) \phi(x_2)$, давайте возьмем $\lambda=\frac{1}{2}$ \\
  И пусть $z := \frac{x_1}{2} + \frac{x_2}{2}$, тогда
  \[
  \phi(z) \le \frac{1}{2} (\phi(x_1) + \phi(x_2))
  \]
  Теперь можно раскрыть значение $\phi$,
  \[
  f(z) - \frac{\mu}{2} \| z \|^2 \le \frac{1}{2} \left(  f(x_1) + f(x_2) - \frac{\mu}{2} (\| x_1\|^2 + \| x_2 \|^2)  \right)
  \]
  Но мы знаем, что $f(x_1) = f(x_2)$, можно записать их сумму, как $2m$.
  Рассмотрим $\|z\|^2$
  \begin{multline*}
  \| z\|^2 = \left\langle z; z \right\rangle = \left\langle \frac{1}{2} (x_1 + x_2); \frac{1}{2} (x_1 + x_2) \right\rangle = \\
  \frac{1}{4} \left\langle x_1 + x_2; x_1 + x_2 \right\rangle =
  \frac{1}{4} \left( \| x_1 \|^2 + 2 \left\langle x_1; x_2 \right\rangle +\| x_2 \|^2 \right)
  \end{multline*}
  Давайте подставим все наши преобразования в неравенство
  \[
  f(z) - \frac{\mu}{8} (\| x_1 \|^2 + 2 \left\langle x_1; x_2 \right\rangle +\| x_2 \|^2) \le m - \frac{\mu}{4} (\|x_1\|^2 + \|x_2\|^2)
  \]
  Давайте теперь сгруппируем члены с $\mu$
  \[
  f(z) \le m - \frac{\mu}{8} (2 \| x_1\|^2 + 2\|x_2\|^2 - \|x_1\|^2 - 2 \left\langle x_1;x_2 \right\rangle - \| x_2 \|^2) = m - \frac{\mu}{8} \| x_1 - x_2 \|^{2}
  \]

  При этом $\mu > 0$, а еще $x_1 \ne x_2$, следовательно их разность нулевой быть не может, тогда $\| x_1 - x_2 \|^2 > 0$ \\
  Из этого получаем, что
  \[
  f(z) \le m - \frac{\mu}{8} \| x_1 - x_2 \|^{2} < m
  \]
  Получили противоречие тому, что $m$~--- минимальное значение.
\end{proof}
  % Единственность минимума сильно выпуклой функции -- нет в конспекте

\InputIfFileExists{text/94DADA28-87E4-4E36-B289-8415BCE750FC-before}{}{}
% generated by roam.py handbook export -- не редактировать, править заметки
% 20260922153328-критерии_сильнои_выпуклости_в_терминах_вторых_производных.org (body 4bfb31b484f3)

\begin{theorem}[{Критерий сильной выпуклости в терминах вторых производных}]\label[theorem]{note:94DADA28-87E4-4E36-B289-8415BCE750FC}
  Пусть $E \subset \R^n$~--- открытое \hyperref[note:316B5BFA-8B9F-4267-9086-D6593AAB5F9B]{выпуклое множество}, $f: E \rightarrow \R$, $f \in C^2$ (дважды \hyperref[note:957FA09D-8883-40A5-B4C7-237D9F815A42]{дифференцируема}), тогда
  $f$~--- $\mu$-\hyperref[note:22E546FC-186F-4A67-85AA-2C85DCEF9021]{сильно выпукла} $\Leftrightarrow \forall x \in E\ \nabla^2 f(x) \ge \mu I$, где $I$~--- единичная матрица \\
  Также равносильно:

  \begin{itemize}
    \item $\nabla^2 f(x) - \mu I \ge 0\ \forall x \in E$
    \item $\inf_{x \in E} (\lambda_{\mathrm{min}} \nabla^2 f(x)) \ge \mu$
  \end{itemize}
\end{theorem}

\begin{proof}
  Возьмем $\phi(x) = f(x) - \frac{\mu}{2} \|x\|^2$. Она дважды \hyperref[note:957FA09D-8883-40A5-B4C7-237D9F815A42]{дифференцируема}.
  Нужно доказать \hyperref[note:1CA14F2C-637F-41DE-9E97-72075D19E97F]{выпуклость} $\phi$, применим \hyperref[note:86FCE16B-C82E-4416-851C-0E0AD0FDA382]{критерий вторых производных}.

  $\phi(x)$~--- выпуклая $\Leftrightarrow \forall x \in E\ \nabla^2 \phi(x) \ge 0$. Давайте посмотрим на \hyperref[note:592AD816-1D80-48E6-8399-8524FBC1E8F3]{гессиан} $\phi$.
  По линейности $\nabla^2 \phi(x) = \nabla^2 \left( f(x) - \frac{\mu}{2} \| x \|^2 \right) = \nabla^2 f(x)- \mu I$.
  Мы знаем, что $\nabla \left( \frac{\mu}{2} \| x \|^2 \right) = \mu x$, а $\nabla (\mu x) = \mu \nabla (x) = \mu I$.
  Итого мы требуем $\forall x \in E\ \nabla^2\phi(x) = \nabla^2 f(x) - \mu I \ge 0$. Что и требуется условием

  \proofpart{Утверждение про $\lambda_{\min}$}
  \proofstep{$\Rightarrow$}
  Пусть $h \ne 0$~--- \hyperref[note:57D369D6-D748-41A2-9BDA-38D2FFC6601A]{собственный вектор} с числом $\lambda$ для $\nabla^2 f(x)$, где $x$~--- произвольный из $E$ \\
  Еще мы знаем, что $\nabla^2 f(x) -\mu I \ge 0 \Rightarrow \left\langle (\nabla^2 f(x) - \mu I)h; h\right\rangle \ge 0$

  Заметим, что $(\nabla^2 f(x) - \mu I) h = \nabla^2 f(x) h - \mu I h$ \\
  А так как $h$~--- \hyperref[note:57D369D6-D748-41A2-9BDA-38D2FFC6601A]{собственный вектор} $\nabla^2 f(x)$, то это равно $\lambda h - \mu h = (\lambda - \mu) h$ \\
  Итого мы получили $\left\langle (\lambda-\mu) h; h \right\rangle \ge 0 \Leftrightarrow (\lambda - \mu) \| h \|^2 \ge 0$ \\
  Так как $\| h \|^2 > 0$, то $(\lambda - \mu) \| h \|^2 \ge 0$ равносильно $(\lambda - \mu) \ge 0$ \\
  Чтобы $\lambda -\mu \ge 0$ было верно для любого $\lambda$, мы как раз и требуем, чтобы $\inf_x (\lambda_{\mathrm{min}} \nabla^2f(x)) \ge \mu$ ( тогда все собственные числа будут не меньше $\mu$).

  \proofstep{$\Leftarrow$}
  Пусть $\inf_{x \in E} (\lambda_{\mathrm{min}} \nabla^2 f(x)) \ge \mu$, то есть для каждого $x \in E$ все \hyperref[note:9D2905FC-8E68-4E6F-9243-364269F795B0]{собственные значения} $A := \nabla^2 f(x)$ не меньше $\mu$.
  Нужно проверить $\left\langle (A - \mu I)h; h \right\rangle \ge 0$ для \emph{любого} $h \in \R^n$, а не только для собственного вектора.

  $f \in C^2$, поэтому $A$~--- симметричная матрица (\hyperref[note:592AD816-1D80-48E6-8399-8524FBC1E8F3]{гессиан}). По \hyperref[note:3BE0BBCE-AAB0-4CB4-8FA6-45277C0889CF]{спектральной теореме} у нее есть ортонормированный базис из собственных векторов $q_1, \dots, q_n$:
  $A q_i = \lambda_i q_i$, $\left\langle q_i; q_j \right\rangle = 0$ при $i \ne j$, $\| q_i \| = 1$.

  Разложим $h = \sum\limits_{i=1}^n c_i q_i$. Тогда
  \[
  (A - \mu I) h = \sum\limits_{i=1}^n c_i (A q_i - \mu q_i) = \sum\limits_{i=1}^n c_i (\lambda_i - \mu) q_i
  \]
  и по ортонормированности базиса
  \[
  \left\langle (A - \mu I)h; h \right\rangle = \sum\limits_{i=1}^n \sum\limits_{j=1}^n c_i c_j (\lambda_i - \mu) \left\langle q_i; q_j \right\rangle = \sum\limits_{i=1}^n (\lambda_i - \mu) c_i^2 \ge 0,
  \]
  так как каждое $\lambda_i \ge \mu$. Значит, $\nabla^2 f(x) - \mu I \ge 0$ для всех $x \in E$.

  \begin{remark}
    Здесь существенна симметричность. У несимметричной матрицы собственные значения ничего не говорят о знаке $\left\langle Ah; h \right\rangle$: у $A = \begin{pmatrix} 1 & -4 \\ 0 & 1 \end{pmatrix}$ оба собственных значения равны $1$, но при $h = (1, 1)^T$ получаем $\left\langle Ah; h \right\rangle = -2 < 0$.
  \end{remark}
\end{proof}

\begin{example}
  Пусть $f(x) = \frac{1}{2} x^T A x$, $A = A^T$. Тогда $\nabla f(x) = Ax$ и $\nabla^2 f(x) = A$ для всех $x$, поэтому $f$~--- $\mu$-сильно выпукла ровно при $0 < \mu \le \lambda_{\mathrm{min}}(A)$.

  \begin{itemize}
    \item $A = \mathrm{diag}(2, 1)$: $f(x) = x_1^2 + \frac{1}{2} x_2^2$~--- $1$-сильно выпукла.
    \item $A = \mathrm{diag}(1, 0)$: $f(x) = \frac{1}{2} x_1^2$ выпукла, но не сильно выпукла. Вдоль оси $x_2$ она постоянна, поэтому у нее бесконечно много точек минимума~--- ср. с \hyperref[note:51CC880D-BF19-43E2-9599-FB31A9A7367F]{единственностью минимума сильно выпуклой функции}.
  \end{itemize}
\end{example}
  % Критерий сильной выпуклости в терминах вторых производных -- теорема 2.2 (в конспекте «дифференциальный критерий»)

\section{L-гладкие функции}\label{lecture:2:2}  % конспект 2.2
\InputIfFileExists{text/02-2-intro}{}{}

\InputIfFileExists{text/7E287F4B-AC3E-4CF5-9D32-3BC2BC52D01A-before}{}{}
% generated by roam.py handbook export -- не редактировать, править заметки
% 20260922164226-l_гладкие_функции.org (body 20f2e75c3d8b)

\begin{definition}[{L-гладкие функции}]\label[definition]{note:7E287F4B-AC3E-4CF5-9D32-3BC2BC52D01A}
  Пусть $f: E \rightarrow \R$, $f \in C^n$ ($n$ раз \hyperref[note:957FA09D-8883-40A5-B4C7-237D9F815A42]{дифференцируема}), $L > 0$, $k \le n$.
  Тогда мы говорим, что  $f$  имеет $L$-гладкие $k$-е производные, если
  \[
  \forall x,y \in E\ \| \nabla^k f(x) - \nabla^k f(y) \|_2 \le L \|x - y\|_2
  \]
  Альтернативное обозначение: $f \in C^{n,k}_L$
\end{definition}

\begin{remark}
  В частности при $k=1$, $f$~--- \hyperref[note:7DD1C7AB-1A49-4675-9F21-5A16C717F670]{функция с липшицевым градиентом} ($f \in C^{n,1}_L$)
\end{remark}

\begin{remark}
  Норма берется по типу $\nabla^k f$: для $k = 1$ (градиент~--- вектор)~--- евклидова, для $k = 2$ (гессиан~--- матрица)~--- спектральная. Для $f \in C^2$ условие при $k = 1$ равносильно $\| \nabla^2 f(x) \| \le L$: см. \hyperref[note:4B6D53BA-E599-4DEA-A557-550BD32042BA]{связь спектральной нормы гессиана и липшицевости градиента} и \hyperref[note:DAA0DAE4-4B48-4A2F-9459-54790B6B129F]{оценку нормы гессиана сверху}.
\end{remark}
  % L-гладкие функции -- опр. 2.2

\InputIfFileExists{text/7DD1C7AB-1A49-4675-9F21-5A16C717F670-before}{}{}
% generated by roam.py handbook export -- не редактировать, править заметки
% 20260922165110-функции_с_липшицевым_градиентом.org (body 581ec3dbcfaf)

\begin{definition}[{Функции с липшицевым градиентом}]\label[definition]{note:7DD1C7AB-1A49-4675-9F21-5A16C717F670}
  Это функции $f: E \rightarrow \R$, имеющие первые $L$-\hyperref[note:7E287F4B-AC3E-4CF5-9D32-3BC2BC52D01A]{гладкие производные} или $f \in C_L^{n,1}$ \\
  Можно записать так (\hyperref[note:1756907B-1FF2-4F62-8553-759B8509281C]{градиент}~--- $L$-липшицево отображение):
  \[
  \| \nabla f(x) - \nabla f(y)\|_2 \le L \| x - y\|_2\ \forall x,y \in E
  \]
\end{definition}

\begin{example}
  \begin{itemize}
    \item $f(x) = \frac{1}{2} x^T A x$, $A = A^T$: $\nabla f(x) = Ax$ и $\| \nabla f(x) - \nabla f(y) \| = \| A(x-y) \| \le \| A \|_2 \| x - y \|$, то есть $L = \| A \|_2$.
    \item $f(x) = x^2$: градиент $2x$ липшицев с $L = 2$, но сама $f$ не липшицева: $\abs{f(x) - f(0)} = x^2$ растет быстрее любой $L \abs{x}$. Поэтому «липшицева функция» и «функция с липшицевым градиентом»~--- разные понятия.
  \end{itemize}
\end{example}
  % Функции с липшицевым градиентом -- опр. 2.2, случай $k = 1$

\InputIfFileExists{text/379E4CAD-A44B-4631-8DAF-F3DC708EAEDE-before}{}{}
% generated by roam.py handbook export -- не редактировать, править заметки
% 20260922165613-усиленное_неравенство_для_функции_с_липшицевым_градиентом.org (body 747b6a33a7b7)

\begin{theorem}[{Усиленное неравенство для функций с липшицевым градиентом}]\label[theorem]{note:379E4CAD-A44B-4631-8DAF-F3DC708EAEDE}
  Пусть $E \subset \R^n$~--- \hyperref[note:316B5BFA-8B9F-4267-9086-D6593AAB5F9B]{выпуклое множество}, $f: E \rightarrow \R$, $f \in C^{1,1}_L$~--- функция с \hyperref[note:7DD1C7AB-1A49-4675-9F21-5A16C717F670]{липшицевым градиентом} \\
  Тогда
  \[
  \forall x,y \in E\ \left| f(y) -f(x) - \left\langle \nabla f(x); y-x  \right\rangle \right|\le \frac{L}{2} \| y - x\|^2_2
  \]
\end{theorem}

\begin{remark}

  Если мы раскроем модуль и посмотрим на верхнюю границу, то получим:
  \[
  \forall x,y \in E\ f(y) - f(x) - \left\langle \nabla f(x); y-x \right\rangle \le \frac{L}{2} \| y -x \|^2
   \]
   Что равносильно $f(y) \le f(x) + \left\langle \nabla f(x); y-x \right\rangle + \frac{L}{2} \| y-x \|^2$,
   а это означает, что график функции лежит не выше параболы~--- двойственно \hyperref[note:22E546FC-186F-4A67-85AA-2C85DCEF9021]{сильной выпуклости}, где он лежит не ниже параболы
\end{remark}

\begin{proof}
  Заметим, что для функции $\phi \in C^1[0;1]$ верно (\hyperref[note:AD362344-ECDC-4627-B344-8356744E6E53]{формула Ньютона-Лейбница}), что $\phi(1) - \phi(0) = \int\limits_0^1 \phi^{\prime}(x)dx$ \\
  Давайте рассмотрим $\phi(t) = f(x + t(y-x))$, тогда $\phi(1) = f(y),\ \phi(0)=f(x)$ \\
  Тогда давайте представим $f(y)-f(x)$, как интеграл  от $\phi^{\prime}$.
  Мы знаем, что $\phi^{\prime}(t) = f^{\prime}(x + t(y-x)) (y-x) = \left\langle \nabla f(x + t(y-x)); y-x \right\rangle$

  Давайте это применим
  \begin{multline*}
  |f(y) -f(x) - \left\langle \nabla f(x); y-x  \right\rangle | =  \\
  \left| \int\limits_0^1 \left\langle \nabla f(x+t(y-x)); y-x \right\rangle dt - \left\langle \nabla f(x); y-x \right\rangle \right|
  \end{multline*}

  Так как $\left\langle \nabla f(x); y-x \right\rangle$~--- константа, ее можно представить так $\int\limits_0^1 \left\langle \nabla f(x); y-x \right\rangle dt$

  Давайте теперь сгруппируем интегралы, а подынтегральную сумму скалярных произведений
  схлопнем по линейности

  \begin{multline*}
  \left| \int\limits_0^1 \left\langle \nabla f(x+t(y-x)); y-x \right\rangle dt - \left\langle \nabla f(x); y-x \right\rangle \right| = \\
  \left| \int\limits_0^1 \left\langle \nabla f(x + t(y-x)) - \nabla f(x); y-x \right\rangle dt\right|
  \end{multline*}

  Еще мы знаем, что модуль интеграла не больше интеграла модуля

  \begin{multline*}
  \left| \int\limits_0^1 \left\langle \nabla f(x + t(y-x)) - \nabla f(x); y-x \right\rangle dt\right| \le
   \int\limits_0^1 \left| \left\langle \nabla f(x + t(y-x)) - \nabla f(x); y-x \right\rangle \right| dt
  \end{multline*}

  А также мы можем применить неравенство \hyperref[note:D0601ACD-8C91-4A93-B0E4-AC0D529C6A5C]{КБШ} для скалярного произведения
  \[
  \left| \left\langle \nabla f(x + t(y-x)) - \nabla f(x); y-x \right\rangle \right| \le \| \nabla f(x + t(y-x)) - \nabla f(x) \| \cdot \| y-x \|
  \]

  Если взять \hyperref[note:7DD1C7AB-1A49-4675-9F21-5A16C717F670]{условие липшицевости градиента} для точек $x + t(y-x)$ и $x$, то получим
  \[
  \| \nabla f(x +  t(y-x)) - \nabla f(x) \| \le L \| x + t(y-x) - x \| = t L \| y-x\|
  \]

  Итого
  \[
  \left| f(y) -f(x) - \left\langle \nabla f(x); y-x \right\rangle \right| \le \int\limits_{0}^1 t L \| y -x \|  \| y -x \| dt =  \frac{L}{2} \| y -x \|^2
  \]
\end{proof}
  % Усиленное неравенство для функций с липшицевым градиентом -- теорема 2.3, замечание 2.2 — в /Rem/

\InputIfFileExists{text/3C6CF60E-0773-488E-A0B9-4D7510D456C0-before}{}{}
% generated by roam.py handbook export -- не редактировать, править заметки
% 20260922174919-связь_сильнои_выпуклости_и_липшицевости_градиента.org (body 119302da06f7)

\begin{theorem}[{Связь сильной выпуклости и липшицевости градиента}]\label[theorem]{note:3C6CF60E-0773-488E-A0B9-4D7510D456C0}
  Пусть $f$~--- $\mu$-\hyperref[note:22E546FC-186F-4A67-85AA-2C85DCEF9021]{сильно выпукла} и $f \in C^{1,1}_L$~--- функция с \hyperref[note:7DD1C7AB-1A49-4675-9F21-5A16C717F670]{липшицевым градиентом}.
  Тогда $\mu \le L$
\end{theorem}

\begin{remark}
  Зачем: для градиентного спуска $\| x_k - x^{*} \|^2 \le \left( 1 - \frac{\mu}{L} \right)^k \| x_0 - x^{*} \|^2$, и $\mu \le L$ гарантирует $1 - \frac{\mu}{L} \in [0;1)$. Отношение $\kappa = \frac{L}{\mu} \ge 1$ называют числом обусловленности: чем оно больше, тем медленнее сходимость.
\end{remark}

\begin{remark}
  Интуиция: сильная выпуклость означает, что функция ограничена параболой снизу, а липшицевость градиента~--- что функция ограничена параболой сверху. Тогда $\mu,L$ отвечают за растянутость этих парабол, следовательно парабола снизу должна быть не уже, чем сверху, отсюда и $\mu \le L$.
\end{remark}

\begin{proof}
  Давайте обозначим $r(x;y) := f(y) - f(x) - \left\langle \nabla f(x); y-x \right\rangle$~--- остаток аппроксимации первого порядка.

  Пусть $E$~--- область определения $f$.
  По \hyperref[note:1B5A3E25-56E8-4D5E-9CF6-E2731DA77604]{дифференциальному критерию сильной выпуклости}
  \[
  \forall x,y \in E\ f(y) \ge f(x) + \left\langle \nabla f(x); y -x \right\rangle + \frac{\mu}{2} \| y -x \|^2
  \]
  Давайте перенесем некоторые слагаемые налево, получим
  \[
  \forall x,y \in E\ f(y) - f(x) - \left\langle \nabla f(x); y -x \right\rangle \ge \frac{\mu}{2} \| y -x \|^2
  \Leftrightarrow r(x;y) \ge \frac{\mu}{2} \|y-x\|^2
  \]

  Теперь оценим $r$ сверху через \hyperref[note:379E4CAD-A44B-4631-8DAF-F3DC708EAEDE]{усиленное неравенство для функций с липшицевым градиентом}
  \[
  \forall x,y \in E\ \left| f(y) -f(x) - \left\langle \nabla f(x); y-x  \right\rangle \right|\le \frac{L}{2} \| y - x\|^2_2
  \]
  Это равносильно
  \[
  \forall x,y \in E\ - \frac{L}{2} \| y - x\|^2_2 \le  f(y) -f(x) - \left\langle \nabla f(x); y-x  \right\rangle \le \frac{L}{2} \| y - x\|^2_2
  \]
  Воспользуемся только правой частью, получим
  \[
  r(x;y) \le \frac{L}{2} \| y-x \|^2
  \]
  Итого
  \[
  \frac{\mu}{2} \| y-x \|^2 \le r(x;y) \le \frac{L}{2} \| y-x\|^2 \Rightarrow \frac{\mu}{2} \| y-x\|^2 \le \frac{L}{2} \|y-x\|^2
  \]
  Давайте возьмем $y \ne x$ и разделим обе части на $\frac{1}{2} \| y-x \|^2$.
  Получим $\mu \le L$.
\end{proof}
  % Связь сильной выпуклости и липшицевости градиента -- в конспекте замечание 3.6 (Пара 3)

\InputIfFileExists{text/4B6D53BA-E599-4DEA-A557-550BD32042BA-before}{}{}
% generated by roam.py handbook export -- не редактировать, править заметки
% 20260922175350-связь_спектральнои_нормы_гессиана_и_липшицевости_градиента.org (body 083af92f09d9)

\begin{theorem}[{Связь спектральной нормы гессиана и липшицевости градиента}]\label[theorem]{note:4B6D53BA-E599-4DEA-A557-550BD32042BA}
  Пусть $f: E \rightarrow \R$, $E$~--- \hyperref[note:316B5BFA-8B9F-4267-9086-D6593AAB5F9B]{выпуклое} и $f \in C^2$ \\
  И $\exists L:\ \forall x\ \| \nabla^2 f(x) \| \le L$, где $\| \cdot \|$~--- спектральная норма (\hyperref[note:592AD816-1D80-48E6-8399-8524FBC1E8F3]{гессиан}) \\
  Тогда $f \in C^{1,1}_L$
\end{theorem}

\begin{remark}
  Верно и обратное~--- \hyperref[note:DAA0DAE4-4B48-4A2F-9459-54790B6B129F]{оценка нормы гессиана сверху}, так что для $f \in C^2$: $\nabla f$~--- $L$-липшицев $\Leftrightarrow \| \nabla^2 f(x) \| \le L\ \forall x$.
\end{remark}

\begin{remark}
  $\nabla^2 f(x)$ симметрична, поэтому по \hyperref[note:3BE0BBCE-AAB0-4CB4-8FA6-45277C0889CF]{спектральной теореме} $\| \nabla^2 f(x) \| = \max_i \abs{\lambda_i}$, и условие равносильно $-L I \le \nabla^2 f(x) \le L I$. Для $\mu$-сильно выпуклой $f$ получаем $\mu I \le \nabla^2 f(x) \le L I$, откуда снова видно $\mu \le L$.
\end{remark}

\begin{proof}
  Вспомним \hyperref[note:5821D7FE-1F9F-4578-BA2B-D60FDB0C69D3]{теорему Лагранжа для векторнозначных функций}:
  Пусть $g: E \rightarrow \R^{n}$,  $E$~--- выпукло и $g$~--- дифференцируема на $E$ \\
  Тогда $\forall a,b \in E\ \exists c \in [a;b]:\ \| g(b) - g(a) \| \le \| g^{\prime}(c) \| \cdot \| b -a \|$

  Рассмотрим $g(x) := \nabla f(x)$ и получим
  \[
  \| \nabla f(a) - \nabla f(b) \| \le \| \nabla^2 f(c) \| \cdot \| b - a \|  \le  L \| b - a \|
  \]
  Получили определение \hyperref[note:7DD1C7AB-1A49-4675-9F21-5A16C717F670]{липшицевости градиента}
\end{proof}
  % Связь спектральной нормы гессиана и липшицевости градиента -- следствие 2.1

\InputIfFileExists{text/DAA0DAE4-4B48-4A2F-9459-54790B6B129F-before}{}{}
% generated by roam.py handbook export -- не редактировать, править заметки
% 20260922181524-оценка_нормы_гессиана_сверху_для_функции_с_липшицевым_градиентом.org (body 9d5497de53e8)

\begin{theorem}[{Оценка нормы гессиана сверху для функции с липшицевым градиентом}]\label[theorem]{note:DAA0DAE4-4B48-4A2F-9459-54790B6B129F}
  Пусть $E \subset \R^n$~--- открытое множество, $f: E \rightarrow \R$, $f \in C^{2,1}_L$: $f$ дважды непрерывно дифференцируема и ее \hyperref[note:7DD1C7AB-1A49-4675-9F21-5A16C717F670]{градиент липшицев} с константой $L$.

  Тогда  $\| \nabla^2 f(x) \| \le L\ \forall x \in E$, где $\| \cdot \|$~--- спектральная норма
\end{theorem}

\begin{remark}
  Обратное~--- \hyperref[note:4B6D53BA-E599-4DEA-A557-550BD32042BA]{связь спектральной нормы гессиана и липшицевости градиента}.
\end{remark}

\begin{proof}
  Возьмем $h \in \R^n$ и $t > 0$. $\nabla f$~--- дифференцируема $\Rightarrow \nabla f(x + th) - \nabla f(x) = \nabla^2 f(x)th + o(th)$ \\
  Тогда нормы двух частей тоже равны:
  \[
  \|  \nabla^2 f(x)th + o(th) \| = \| \nabla f(x + th) - \nabla f(x) \| \le L \| x + th - x \|
  \]
  По липшицевости градиента.
  Давайте разделим обе части неравенства на $t$, и устремим $t \rightarrow 0+$
  \[
  \|  \nabla^2 f(x)h + \frac{o(th)}{t}  \| \le L \| h \|
  \]
  Мы знаем, что $\frac{o(th)}{t} \rightarrow 0$. Тогда по предельному переходу
  \[
  \|  \nabla^2 f(x)h  \| \le L \| h \|
  \]
  Так как $h$ произвольный, по определению спектральной нормы $\| \nabla^2 f(x) \| = \sup\limits_{h \ne 0} \frac{\| \nabla^2 f(x) h \|}{\| h \|} \le L$.
\end{proof}
  % Оценка нормы гессиана сверху для функции с липшицевым градиентом -- нет в конспекте (обратное к следствию 2.1)

\InputIfFileExists{text/81C8B68E-DAAD-4B86-BAEB-AB0E1C07442F-before}{}{}
% generated by roam.py handbook export -- не редактировать, править заметки
% 20260922182611-аппроксимация_функции_с_липшицевым_гессианом_до_3_степени.org (body 62a5bf447270)

\begin{theorem}[{Аппроксимация функции с липшицевым гессианом до 3 степени}]\label[theorem]{note:81C8B68E-DAAD-4B86-BAEB-AB0E1C07442F}
  Пусть $f: \R^n \rightarrow \R$, $f \in C^{2,2}_L$ (\hyperref[note:7E287F4B-AC3E-4CF5-9D32-3BC2BC52D01A]{липшицев гессиан}). Тогда
  \[
  \left| f(x+y) - f(x) - \left\langle \nabla f(x) ; y \right\rangle -
  \frac{1}{2} \left\langle \nabla^2f(x) y; y \right\rangle \right| \le \frac{L}{6} \|y\|^{3}
  \]
\end{theorem}

\begin{proof}
  Аналогично \hyperref[note:379E4CAD-A44B-4631-8DAF-F3DC708EAEDE]{усиленному неравенству для функций с липшицевым градиентом}

  Давайте возьмем $\phi(t) = f(x + ty)$. Мы знаем, что (\hyperref[note:AD362344-ECDC-4627-B344-8356744E6E53]{Ньютон-Лейбниц}) $\phi(1) - \phi(0)= \int\limits_0^1 \phi^{\prime}(t) dt$.
  Но нам нужно подняться на порядок выше, чем в \hyperref[note:379E4CAD-A44B-4631-8DAF-F3DC708EAEDE]{усиленном неравенстве}.
  Давайте тогда выразим $\phi^{\prime}(t)$ также через Ньютона-Лейбница (можем сделать так, так как $\phi^{\prime\prime}$ непрерывна: $f \in C^2$). $\phi^{\prime}(t) - \phi^{\prime}(0) = \int\limits_0^t \phi^{\prime\prime}(u) du$.  (пришлось еще добавить $\phi^{\prime}(0)$, чтобы получилась разность).
  Тогда $\phi(1) - \phi(0) = \int\limits_0^1 \int\limits_0^t \phi^{\prime\prime}(u) du dt + \phi^{\prime}(0)\Rightarrow \int\limits_0^1 \int\limits_0^t \phi^{\prime\prime}(u) du dt = \phi(1) - \phi(0) - \phi^{\prime}(0)$ \\
  Давайте теперь посчитаем сами производные
  $\phi^{\prime}(t) = \left\langle \nabla f(x + ty); y \right\rangle$, $\phi^{\prime\prime}(t) = \left\langle \nabla^2 f(x + ty)y; y \right\rangle$ (второе слагаемое сократилось так как $y^{\prime}_t = 0$).

  Заметим, что $\phi(1) = f(x+y), \phi(0) = f(x), \phi^{\prime}(0) = \left\langle \nabla f(x); y \right\rangle, \phi^{\prime\prime}(0) = \left\langle \nabla^2 f(x)y; y \right\rangle$.
  Подставим это в модуль из неравенства  и получим $\left| \phi(1) - \phi(0) - \phi^{\prime}(0) - \frac{1}{2} \phi^{\prime\prime} (0) \right|$.

  Заметим, что первая тройка заменяется интегралом, а так как $\frac{1}{2} \phi^{\prime\prime}(0)$~--- константа, ее тоже можно занести под интеграл: $c \cdot \int\limits_0^1 \int\limits_0^t du dt  = c \cdot \int\limits_0^1 t dt =  c \cdot \frac{1}{2} \Rightarrow \frac{1}{2}\phi^{\prime\prime}(0) = \int\limits_0^1 \int\limits_0^t \phi^{\prime\prime}(0) du dt$

  Давайте запишем модуль, как объединенный интеграл
  \[
  \left| \phi(1) - \phi(0) - \phi^{\prime}(0) - \frac{1}{2} \phi^{\prime\prime} (0) \right| = \left| \int\limits_0^1 \int\limits_0^t \left[ \phi^{\prime\prime}(u)-\phi^{\prime\prime}(0) \right]du dt  \right|
  \]
  Оценим интегралом от модуля сверху
  \[
  \left| \int\limits_0^1 \int\limits_0^t \left[ \phi^{\prime\prime}(u)-\phi^{\prime\prime}(0) \right]du dt  \right| \le
   \int\limits_0^1 \int\limits_0^t \left| \phi^{\prime\prime}(u)-\phi^{\prime\prime}(0) \right|du dt
  \]
  Давайте теперь подставим вторые производные
  \[
  \int\limits_0^1 \int\limits_0^t \left| \phi^{\prime\prime}(u)-\phi^{\prime\prime}(0) \right|du dt =
   \int\limits_0^1 \int\limits_0^t \left| \left\langle \nabla^2 f(x + uy)y; y  \right\rangle - \left\langle \nabla^2f(x)y;  y \right\rangle \right|du dt
  \]
  Свернем скалярные произведения по аргументу
  \[
   \int\limits_0^1 \int\limits_0^t \left| \left\langle \nabla^2 f(x + uy)y; y  \right\rangle - \left\langle \nabla^2f(x)y;  y \right\rangle \right|du dt  =
   \int\limits_0^1 \int\limits_0^t \left| \left\langle \nabla^2 f(x + uy)y - \nabla^2 f(x) y; y  \right\rangle \right|du dt
  \]
  Применим \hyperref[note:D0601ACD-8C91-4A93-B0E4-AC0D529C6A5C]{КБШ} для скалярного произведения
  \[
   \int\limits_0^1 \int\limits_0^t \left| \left\langle \nabla^2 f(x + uy)y - \nabla^2 f(x) y; y  \right\rangle \right|du dt \le
   \int\limits_0^1 \int\limits_0^t \| \nabla^2 f(x + uy)y - \nabla^2f(x)y\| \cdot \| y\| du dt
  \]

  Заметим, что $\| \nabla^2 f(x + uy)y - \nabla^2f(x)y\| = \| (\nabla^2 f(x + uy) - \nabla^2f(x)) y\|$ \\
  Если взять матрицу $A:= (\nabla^2 f(x + uy) - \nabla^2f(x))$, то
  $\| A \| = \sup\limits_{h \ne 0} \frac{\| A h \|}{\| h\|} \Rightarrow \|Ah \| \le \|A\| \|h\|\ \forall h$ \\
  Тогда оценим сверху $\|Ay\|$
  \[
   \int\limits_0^1 \int\limits_0^t \| \nabla^2 f(x + uy)y - \nabla^2f(x)y\| \cdot \| y\| du dt  \le
   \int\limits_0^1 \int\limits_0^t \| \nabla^2 f(x + uy) - \nabla^2f(x)\| \cdot \| y\|^2 du dt
  \]
  Тут у нас разность гессианов, так как у нас $L$ гладкие 2 производные, то можем еще оценить сверху
  \[
   \int\limits_0^1 \int\limits_0^t \| \nabla^2 f(x + uy) - \nabla^2f(x)\| \cdot \| y\|^2 du dt
   \le
   \int\limits_0^1 \int\limits_0^t L \| x + uy - x \| \cdot \| y\|^2 du dt =
   \int\limits_0^1 \int\limits_0^t L u \| y \| \cdot \| y\|^2 du dt
  \]
  Давайте теперь вынесем константы за интеграл и посчитаем
  \[
   \int\limits_0^1 \int\limits_0^t L u \| y \| \cdot \| y\|^2 du dt =
   L \|y\|^3  \cdot \int\limits_0^1 \int\limits_0^t u\, du dt = L \|y\|^3 \cdot \int\limits_0^1 \frac{t^2}{2}\,dt
   = \frac{L}{6} \| y \|^3
  \]

  Итого левая часть неравенства не больше $\frac{L}{6} \| y \|^3$.
\end{proof}
  % Аппроксимация функции с липшицевым гессианом до 3 степени -- теорема 2.4

\InputIfFileExists{text/655064DF-D686-4726-83DC-6F3D67CC5B9C-before}{}{}
% generated by roam.py handbook export -- не редактировать, править заметки
% 20260922182947-оценка_для_выпуклых_функции_с_липшицевым_градиентом.org (body 8829a1d674d5)

\begin{theorem}[{Оценка для выпуклых функций с липшицевым градиентом}]\label[theorem]{note:655064DF-D686-4726-83DC-6F3D67CC5B9C}
  Пусть $f: \R^n \rightarrow \R$, $f \in C^{1,1}_L$ (\hyperref[note:7DD1C7AB-1A49-4675-9F21-5A16C717F670]{липшицев градиент}) и \hyperref[note:1CA14F2C-637F-41DE-9E97-72075D19E97F]{выпуклая};
  Тогда
  \[
  \forall x, y \in \R^n\ f(x) + \left\langle \nabla f(x); y-x \right\rangle + \frac{1}{2L} \| \nabla f(y) - \nabla f(x) \|^2 \le f(y)
  \]
\end{theorem}

\begin{proof}
  Давайте рассмотрим $\phi(y) := f(y) + \left\langle - \nabla f(x); y \right\rangle$.
  Заметим, что $\phi$ выпуклая, так как $f$~--- выпуклая, а скалярное произведение~--- линейная функция;
  Тогда $\phi$~--- выпуклая, как сумма выпуклой и линейной функций

  Заметим, что у $\phi$ тоже липшицев градиент.
  \begin{multline*}
  \nabla \phi (y) = \nabla f (y) - \nabla f(x) \Rightarrow \| \nabla \phi(y_1) - \nabla \phi(y_2) \| = \\
  \| \nabla f (y_1) - \nabla f(x) - \nabla f (y_2) + \nabla f(x) \| = \| \nabla f(y_1) - \nabla f(y_2) \|
  \end{multline*}
  А мы знаем, что у $f$ липшицев градиент. Тогда
  \[
  \| \nabla \phi(y_1) - \nabla \phi(y_2) \| = \| \nabla f(y_1) - \nabla f(y_2) \| \le L \|y_1 - y_2 \|
  \]

  Заметим, что $\nabla \phi(x) = \nabla f(x) - \nabla f(x) = 0$. $\phi$ выпуклая, поэтому по \hyperref[note:035B8935-65A9-4BB4-BDA3-60A60FF278F0]{теореме о нулевом градиенте} $x$~--- точка наименьшего значения $\phi$.

  Тогда $\phi(x) \le \phi(z)\ \forall z \in \R^n$.
  Возьмем $z = y - \frac{1}{L} \nabla \phi(y)$.

  По \hyperref[note:379E4CAD-A44B-4631-8DAF-F3DC708EAEDE]{усиленному неравенству} для $\phi$ (у нее липшицев градиент) $\phi(z)-\phi(y) - \left\langle \nabla \phi(y);z - y \right\rangle \le \frac{L}{2} \| z -y \|^2$ \\
  Давайте оценим $\phi(z)$. И подставим $z$ в правой части
  \[
  \phi(z) \le \phi(y) +  \left\langle \nabla \phi (y); y - \frac{1}{L} \nabla \phi (y) - y \right\rangle + \frac{L}{2} \left\| y - \frac{1}{L} \nabla \phi (y) - y \right\|^{2}
  \]
  Сократим
  \[
  \phi (z) \le  \phi(y) - \frac{1}{L} \left\langle \nabla \phi(y); \nabla \phi(y) \right\rangle + \frac{1}{2L} \| \nabla \phi(y) \|^2 =
  \phi(y) - \frac{1}{2L} \| \nabla \phi (y) \|^2
  \]

  Получили $\phi(x) \le \phi(y) - \frac{1}{2L} \| \nabla \phi (y) \|^2$, подставим значение $\phi$
  \[
  f(x) - \left\langle \nabla f(x); x \right\rangle \le f(y) - \left\langle \nabla f(x);y \right\rangle - \frac{1}{2L} \| \nabla f(y) - \nabla f(x) \|^2
  \]
  Если оставить справа только $f(y)$ и сократить скалярные, то получим
  \[
  f(x) + \left\langle \nabla f(x); y -x \right\rangle + \frac{1}{2L} \| \nabla f(y) - \nabla f(x) \|^2  \le f(y)
  \]
\end{proof}
  % Оценка для выпуклых функций с липшицевым градиентом -- теорема 2.5

\InputIfFileExists{text/02-outro}{}{}
